\subsection{Complex and Real Jordan Forms and Block Sizes}

Jordan bases establish existence; shifted-power ranks recover block multiplicities and uniqueness.

\begin{theorem} \label{thm:complex-jordan-form} \hlt{Jordan Form over $\mathbb{C}$}\\
Let $\mathcal{V}$ be a complex vector space with $\dim_{\mathbb{C}}\mathcal{V}=n\in\mathbb{Z}_{>0}$, and let $T\in\LL(\mathcal{V})$.\\
Then $\exists$ an ordered basis $\mathcal{B}$ of $\mathcal{V}$ and a Jordan matrix $J$ such that $[T]_{\mathcal{B}\leftarrow\mathcal{B}}=J$.
\begin{proof}
\begin{enumerate}[label=(\roman*)]
\setlength{\itemsep}{0pt}
\item \hlt{Primary Splitting}: Let $\operatorname{spec}(T)=\{\lambda_1,\ldots,\lambda_m\}$ and set $G_j\defeq\mathcal{G}_{\lambda_j}(T)$ for $j\in\{1,\ldots,m\}$. Set $N_j\defeq(T-\lambda_jI_{\mathcal{V}})|_{G_j}$.\\
The identity $\mathcal{V}=\bigoplus\limits_{j=1}^{m}G_j$ holds, and every $N_j$ is nilpotent, by Theorem \ref{thm:complex-primary-decomposition-generalized-eigenspaces}(i)--(ii).
\item \hlt{Chain Bases}: For $j\in\{1,\ldots,m\}$, choose an $N_j$-chain basis $\mathcal{B}_j$ of $G_j$ by Theorem \ref{thm:nilpotent-jordan-chain-basis}.
\item \hlt{Block Identification}: On a chain of length $r$ in $\mathcal{B}_j$, the matrix of $N_j$ is $J_r(0)$ by the chain recurrence.\\
The matrix of $T|_{G_j}=\lambda_jI_{G_j}+N_j$ on the same chain is $J_r(\lambda_j)$ by Definition \ref{def:jordan-chain-block-and-matrix}(ii).
\item \hlt{Basis Assembly}: The concatenation $\mathcal{B}\defeq(\mathcal{B}_1,\ldots,\mathcal{B}_m)$ is an ordered basis of $\mathcal{V}$ by the primary direct sum.\\
Matrix $[T]_{\mathcal{B}\leftarrow\mathcal{B}}$ is matching direct sum of Jordan blocks by item (iii). Hence it is Jordan matrix by Definition \ref{def:jordan-chain-block-and-matrix}(iii).
\end{enumerate}
\hlt{Structural Commentary}: Primary splitting isolates nilpotent shifts; chain bases concatenate into a Jordan basis.\\
\hlt{Key Insight}: Each primary restriction is $\lambda_jI_{G_j}$ plus its nilpotent shift.
\end{proof}
\end{theorem}

Figure \ref{fig:jordan-form-primary-chains} tracks primary summands through nilpotent chains to Jordan blocks.

\begin{illustrationfigure}[H]
  \centering
  \begin{tikzpicture}[
  font=\small,
  block/.style={draw=blue!65!black, fill=blue!6, rounded corners=2pt, minimum width=2.15cm, minimum height=0.68cm, align=center},
  arrow/.style={-{Stealth[length=2mm]}, thick, gray!75}
]
  \node[block] (v) at (0,0) {$\mathcal{V}$};
  \node[block] (g1) at (3.0,0.85) {$\mathcal{G}_{\lambda_1}(T)$};
  \node[block] (gj) at (3.0,0) {$\cdots$};
  \node[block] (gm) at (3.0,-0.85) {$\mathcal{G}_{\lambda_m}(T)$};

  \node[block] (j1) at (6.0,0.85) {$J_k(\lambda_1)$};
  \node[block] (jj) at (6.0,0) {$\cdots$};
  \node[block] (jm) at (6.0,-0.85) {$J_k(\lambda_m)$};

  \draw[arrow] (v.east) -- (g1.west);
  \draw[arrow] (v.east) -- (gj.west);
  \draw[arrow] (v.east) -- (gm.west);
  \draw[arrow] (g1.east) -- (j1.west);
  \draw[arrow] (gj.east) -- (jj.west);
  \draw[arrow] (gm.east) -- (jm.west);
\end{tikzpicture}

  \caption{Primary pieces $\mathcal{G}_{\lambda_j}(T)$ decompose into nilpotent chains, which determine Jordan blocks.}
  \label{fig:jordan-form-primary-chains}
\end{illustrationfigure}

\begin{theorem} \label{thm:real-jordan-block-form} \hlt{Real Jordan Block Form}\\
Let $n\in\mathbb{Z}_{>0}$, $A\in\operatorname{Mat}(n,n,\mathbb{R})$. For $a\in\mathbb{R}$, $b>0$, set $R(a,b)\defeq\left(\begin{smallmatrix}a&b\\-b&a\end{smallmatrix}\right)$. For $k\in\mathbb{Z}_{>0}$, define $\mathcal{J}_k(a,b)\defeq[C_{rs}]_{r,s=1}^{k}\in\operatorname{Mat}(2k,2k,\mathbb{R})$.\\
Its blocks satisfy $C_{rr}=R(a,b)$, $C_{r,r+1}=I_2$ for $r<k$, and $C_{rs}=0_2$ otherwise. Then:
\begin{enumerate}[label=(\roman*)]
\setlength{\itemsep}{0pt}
\item $\exists S\in\operatorname{GL}_n(\mathbb{R})$ such that $S^{-1}AS$ is a direct sum. Its scalar summands are $J_k(\lambda)$ with $(k,\lambda)\in\mathbb{Z}_{>0}\times\mathbb{R}$.\\
Its remaining summands are $\mathcal{J}_k(a,b)$ with $(k,a,b)\in\mathbb{Z}_{>0}\times\mathbb{R}\times\mathbb{R}_{>0}$.
\item For $k\in\mathbb{Z}_{>0}$, $a\in\mathbb{R}$, and $b>0$, $\exists Q\in\operatorname{GL}_{2k}(\mathbb{C})$ such that $Q^{-1}\mathcal{J}_k(a,b)Q=J_k(a+ib)\oplus J_k(a-ib)$.
\end{enumerate}
\begin{proof}
\begin{enumerate}[label=(\roman*)]
\setlength{\itemsep}{0pt}
\item[$(i)$]
\begin{enumerate}[label=(\alph*)]
\setlength{\itemsep}{0pt}
\item \hlt{Primary Partition}: Regard $A$ as an operator on $\mathbb{R}^n$ and let $A_{\mathbb{C}}$ be its complexification.\\
Set $\Lambda_{\mathbb{R}}\defeq\operatorname{spec}(A_{\mathbb{C}})\cap\mathbb{R}$, $\Lambda_+\defeq\{\lambda\in\operatorname{spec}(A_{\mathbb{C}}):\operatorname{Im}(\lambda)>0\}$. For $\mu\in\operatorname{spec}(A_{\mathbb{C}})$, set $G_{\mu}\defeq\mathcal{G}_{\mu}(A_{\mathbb{C}})$. Set $\overline{\Lambda_+}\defeq\{\overline{\lambda}:\lambda\in\Lambda_+\}$.\\
The spectrum is the disjoint union of $\Lambda_{\mathbb{R}}$, $\Lambda_+$, and $\overline{\Lambda_+}$ by Corollary \ref{cor:complexification-conjugate-eigenvalue-multiplicities}(i).\\
Hence $\mathbb{C}^n=\bigoplus\limits_{\lambda\in\Lambda_{\mathbb{R}}}G_{\lambda}\oplus\bigoplus\limits_{\lambda\in\Lambda_+}(G_{\lambda}\oplus G_{\overline{\lambda}})$ by Theorem \ref{thm:complex-primary-decomposition-generalized-eigenspaces}(i).
\item \hlt{Real Primary Chains}: Fix $\lambda\in\Lambda_{\mathbb{R}}$ and set $G_{\lambda}^{\mathbb{R}}\defeq\ker((A-\lambda I_n)^n)$.\\
The equality $(G_{\lambda}^{\mathbb{R}})_{\mathbb{C}}=\ker((A_{\mathbb{C}}-\lambda I_n)^n)$ holds by Proposition \ref{prop:complexification-operator-basic-properties}(ii) and equality of real and imaginary parts.\\
Thus $(G_{\lambda}^{\mathbb{R}})_{\mathbb{C}}=G_{\lambda}$ by Theorem \ref{thm:generalized-eigenspace-stabilized-kernel}(i). Subspace $G_{\lambda}^{\mathbb{R}}$ is invariant under $A-\lambda I_n$ by commutation with its $n$th power.\\
Set $N_{\lambda}\defeq(A-\lambda I_n)|_{G_{\lambda}^{\mathbb{R}}}$; equality $N_{\lambda}^n=0$ holds by definition of $G_{\lambda}^{\mathbb{R}}$. Choose an $N_{\lambda}$-chain basis $\mathcal{B}_{\lambda}$ by Theorem \ref{thm:nilpotent-jordan-chain-basis}.\\
The corresponding matrix blocks of $A|_{G_{\lambda}^{\mathbb{R}}}=\lambda I_{G_{\lambda}^{\mathbb{R}}}+N_{\lambda}$ are $J_k(\lambda)$ by the chain recurrence and Definition \ref{def:jordan-chain-block-and-matrix}.
\item \hlt{Conjugate-Pair Chains}: Fix $\lambda=a+ib\in\Lambda_+$ and set $N_{\lambda}^{\mathbb{C}}\defeq(A_{\mathbb{C}}-\lambda I_n)|_{G_{\lambda}}$.\\
This operator is nilpotent by Theorem \ref{thm:complex-primary-decomposition-generalized-eigenspaces}(ii). Choose an $N_{\lambda}^{\mathbb{C}}$-chain basis by Theorem \ref{thm:nilpotent-jordan-chain-basis}.\\
Each conjugate chain is a $\overline{\lambda}$-chain by conjugating its recurrence and using the real entries of $A$.\\
Conjugation is a conjugate-linear bijection $G_{\lambda}\to G_{\overline{\lambda}}$ by Lemma \ref{lem:complexification-conjugate-generalized-eigenspaces} and its involutivity.\\
Hence the conjugate chains form a basis of $G_{\overline{\lambda}}$ by coefficient conjugation.\\
Fix a chain $\mathcal{C}=(\mathbf{v}_1,\ldots,\mathbf{v}_k)$, set $\mathbf{v}_0\defeq\mathbf{0}$, and write $\mathbf{v}_j=\mathbf{u}_j+i\mathbf{w}_j$ for $j\in\{0,\ldots,k\}$.\\
Set $U_j\defeq[\mathbf{u}_j\ \mathbf{w}_j]$ and $\mathcal{U}_{\mathcal{C}}\defeq(\mathbf{u}_1,\mathbf{w}_1,\ldots,\mathbf{u}_k,\mathbf{w}_k)$.\\
The lists $(\mathcal{C},\overline{\mathcal{C}})$ and $\mathcal{U}_{\mathcal{C}}$ differ by an invertible change of coordinates by $\mathbf{u}_j=\frac{1}{2}(\mathbf{v}_j+\overline{\mathbf{v}}_j)$ and $\mathbf{w}_j=\frac{1}{2i}(\mathbf{v}_j-\overline{\mathbf{v}}_j)$.\\
Its inverse is $\mathbf{v}_j=\mathbf{u}_j+i\mathbf{w}_j$ and $\overline{\mathbf{v}}_j=\mathbf{u}_j-i\mathbf{w}_j$.\\
The identity $AU_j=U_jR(a,b)+U_{j-1}$ holds by separating the real and imaginary parts of the chain recurrence.\\
Thus the matrix of $A$ on $\mathcal{U}_{\mathcal{C}}$ is $\mathcal{J}_k(a,b)$ by the block definition.
\item \hlt{Real Basis Assembly}: Let $\mathcal{B}$ concatenate the bases $\mathcal{B}_{\lambda}$ for $\lambda\in\Lambda_{\mathbb{R}}$ and the lists $\mathcal{U}_{\mathcal{C}}$ over all chosen chains for $\lambda\in\Lambda_+$.\\
The complexified lists $\mathcal{B}_{\lambda}$ are bases of $G_{\lambda}$ by (b) and Proposition \ref{prop:complexification-vector-space-basis}(ii).\\
Hence $\mathcal{B}$ is a basis of $\mathbb{C}^n$ by the primary direct sum and the invertible replacements in (c).\\
The list $\mathcal{B}$ is a real basis of $\mathbb{R}^n$ by its complex linear independence and $\#\mathcal{B}=n$.\\
If $S$ is its basis-column matrix, then $S\in\operatorname{GL}_n(\mathbb{R})$ and $S^{-1}AS$ has the asserted form by the block calculations in (b)--(c).
\end{enumerate}
\item[$(ii)$] Fix $(k,a,b)\in\mathbb{Z}_{>0}\times\mathbb{R}\times\mathbb{R}_{>0}$ and let $(\mathbf{e}_j,\mathbf{f}_j)_{j=1}^{k}$ be the standard coordinate pairs of $\mathbb{C}^{2k}$, with $\mathbf{e}_0=\mathbf{f}_0\defeq\mathbf{0}$.\\
Set $\lambda\defeq a+ib$ and, for $j\in\{1,\ldots,k\}$, set $\mathbf{v}_j\defeq\mathbf{e}_j+i\mathbf{f}_j$.\\
The identity $(\forall j\in\{1,\ldots,k\})\,(\mathcal{J}_k(a,b)-\lambda I_{2k})\mathbf{v}_j=\mathbf{v}_{j-1}$ holds by the block definition.\\
The conjugate identity $(\forall j\in\{1,\ldots,k\})\,(\mathcal{J}_k(a,b)-\overline{\lambda}I_{2k})\overline{\mathbf{v}}_j=\overline{\mathbf{v}}_{j-1}$ holds by entrywise conjugation.\\
The list $(\mathbf{v}_1,\ldots,\mathbf{v}_k,\overline{\mathbf{v}}_1,\ldots,\overline{\mathbf{v}}_k)$ is a basis by the invertible pairwise relations $\mathbf{e}_j=\frac{1}{2}(\mathbf{v}_j+\overline{\mathbf{v}}_j)$ and $\mathbf{f}_j=\frac{1}{2i}(\mathbf{v}_j-\overline{\mathbf{v}}_j)$.\\
If $Q$ is its basis-column matrix, then $Q^{-1}\mathcal{J}_k(a,b)Q=J_k(a+ib)\oplus J_k(a-ib)$ by the two chain recurrences.
\end{enumerate}
\hlt{Structural Commentary}: Real and conjugate primary chains assemble the real blocks and their complex forms.\\
\hlt{Key Insight}: Pairing conjugate chains converts multiplication by $a+ib$ into $R(a,b)$.
\end{proof}
\end{theorem}

\begin{corollary} \label{cor:jordan-blocks-diagonalizability} \hlt{Diagonalizability in Jordan Form}\\
Let $n\in\mathbb{Z}_{>0}$, let $A\in\operatorname{Mat}(n,n,\mathbb{C})$, and let $J=\bigoplus\limits_{\alpha=1}^s J_{k_{\alpha}}(\mu_{\alpha})$ be a Jordan form of $A$.\\
For $(\lambda,\alpha)\in\operatorname{spec}(A)\times\{1,\ldots,s\}$, set $\delta_{\lambda\alpha}\defeq\delta_{\lambda,\mu_{\alpha}}$, and set $b_J(\lambda)\defeq\sum\limits_{\alpha=1}^s\delta_{\lambda\alpha}$. Then:
\begin{enumerate}[label=(\roman*)]
\setlength{\itemsep}{0pt}
\item $(\forall\lambda\in\operatorname{spec}(A))\,\dim E_{\lambda}(A)=b_J(\lambda)$;
\item $A$ is diagonalizable if and only if $(\forall\alpha\in\{1,\ldots,s\})\,k_{\alpha}=1$.
\end{enumerate}
\begin{proof}
Choose $S\in\operatorname{GL}_n(\mathbb{C})$ such that $A=SJS^{-1}$. For $\lambda\in\operatorname{spec}(A)$, set $a_J(\lambda)\defeq\sum\limits_{\alpha=1}^{s}k_{\alpha}\delta_{\lambda\alpha}$.\\
The algebraic multiplicity of $\lambda$ in $A$ equals that in $J$ by Theorem \ref{thm:similarity-preserves-characteristic-polynomial}(ii). The latter multiplicity is $a_J(\lambda)$ by Proposition \ref{prop:upper-triangular-characteristic-polynomial}(iii).
\begin{enumerate}[label=(\roman*)]
\setlength{\itemsep}{0pt}
\item[$(i)$] Fix $\lambda\in\operatorname{spec}(A)$. The map $S$ restricts to an isomorphism $E_{\lambda}(J)\to E_{\lambda}(A)$ by $A=SJS^{-1}$.\\
If $\mu_{\alpha}=\lambda$, then $E_{\lambda}(J_{k_{\alpha}}(\mu_{\alpha}))=\operatorname{span}(\mathbf{e}_1)$ by Definition \ref{def:jordan-chain-block-and-matrix}(ii). If $\mu_{\alpha}\neq\lambda$, then $J_{k_{\alpha}}(\mu_{\alpha})-\lambda I_{k_{\alpha}}$ is invertible by Lemma \ref{lem:scalar-plus-nilpotent-invertible}.\\
Thus $\dim E_{\lambda}(J_{k_{\alpha}}(\mu_{\alpha}))=\delta_{\lambda\alpha}$ and $\dim E_{\lambda}(A)=\sum\limits_{\alpha=1}^{s}\delta_{\lambda\alpha}=b_J(\lambda)$ by the direct-sum block structure.
\item[$(ii)$] ($\Rightarrow$) If $A$ is diagonalizable, then $(\forall\lambda\in\operatorname{spec}(A))\,a_J(\lambda)=b_J(\lambda)$ by Theorem \ref{thm:diagonalizable-iff-algebraic-geometric-multiplicities-match} and part (i).\\
Fix $\alpha\in\{1,\ldots,s\}$ and set $\lambda\defeq\mu_{\alpha}$. Then $k_{\alpha}=1$ by $0=a_J(\lambda)-b_J(\lambda)=\sum\limits_{\beta=1}^{s}(k_{\beta}-1)\delta_{\lambda\beta}$ and nonnegativity of its summands.
\item[$(ii)$] ($\Leftarrow$) Suppose $(\forall\alpha\in\{1,\ldots,s\})\,k_{\alpha}=1$. The matrix $A=SJS^{-1}$ is diagonalizable by diagonality of $J$.
\end{enumerate}
\hlt{Structural Commentary}: Blockwise kernels count eigendirections; multiplicity equality then excludes every longer block.\\
\hlt{Key Insight}: A Jordan block is diagonal exactly when its length is one.
\end{proof}
\end{corollary}

\begin{lemma} \label{lem:nilpotent-jordan-block-power-rank} \hlt{Powers of a Nilpotent Jordan Block}\\
Let $k\in\mathbb{Z}_{>0}$, let $p\in\mathbb{Z}_{\geq0}$, and set $N\defeq J_k(0)\in\operatorname{Mat}(k,k,\mathbb{C})$. Then:
\begin{enumerate}[label=(\roman*)]
\setlength{\itemsep}{0pt}
\item $\operatorname{rank}(N^p)=\max(k-p,0)$;
\item $\dim\ker(N^p)=\min(p,k)$.
\end{enumerate}
\begin{proof}
Let $(\mathbf{e}_j)_{j=1}^{k}$ be the standard basis. Set $(q)_+\defeq\max(q,0)$ and $\mathcal{E}_m\defeq\operatorname{span}(\mathbf{e}_1,\ldots,\mathbf{e}_m)$ for $m\in\{1,\ldots,k\}$, with $\mathcal{E}_0\defeq\{\mathbf{0}\}$.\\
Set $\mathbf{e}_j\defeq\mathbf{0}$ for $j\in\mathbb{Z}_{\leq0}$. Then $(\forall j\in\{1,\ldots,k\})\,N^p\mathbf{e}_j=\mathbf{e}_{j-p}$ by Definition \ref{def:jordan-chain-block-and-matrix}(ii) and induction on $p$.
\begin{enumerate}[label=(\roman*)]
\setlength{\itemsep}{0pt}
\item If $p\geq k$, then $\operatorname{im}(N^p)=\mathcal{E}_0$ by the preceding power identity.\\
If $p<k$, then the nonzero images are exactly $\mathbf{e}_1,\ldots,\mathbf{e}_{k-p}$, so $\operatorname{im}(N^p)=\mathcal{E}_{k-p}$.\\
Thus $\operatorname{im}(N^p)=\mathcal{E}_{(k-p)_+}$ and $\operatorname{rank}(N^p)=(k-p)_+=\max(k-p,0)$ in both cases.
\item The equality $\dim\ker(N^p)=k-(k-p)_+=\min(p,k)$ holds by Corollary \ref{cor:matrix-rank-nullity}(i) and part (i).
\end{enumerate}
\hlt{Structural Commentary}: The block shift identifies the surviving image basis; rank--nullity computes its complement.\\
\hlt{Key Insight}: The $p$th power deletes the first $p$ chain positions.
\end{proof}
\end{lemma}

\begin{definition} \label{def:weyr-characteristic} \hlt{Weyr Characteristic}\\
Let $n\in\mathbb{Z}_{>0}$, let $A\in\operatorname{Mat}(n,n,\mathbb{C})$, and let $\lambda\in\operatorname{spec}(A)$.
\begin{enumerate}[label=(\roman*)]
\setlength{\itemsep}{0pt}
\item For $p\in\mathbb{Z}_{\geq0}$, the \hlt{shifted-power rank} is $r_p(\lambda,A)\defeq\operatorname{rank}((A-\lambda I_n)^p)$.
\item For $p\in\mathbb{Z}_{>0}$, the $p$th \hlt{Weyr number} is $w_p(\lambda,A)\defeq r_{p-1}(\lambda,A)-r_p(\lambda,A)$.
\end{enumerate}
Identity $(\forall p\in\mathbb{Z}_{>0})\,w_p(\lambda,A)=\dim\ker((A-\lambda I_n)^p)-\dim\ker((A-\lambda I_n)^{p-1})\geq0$ holds by Proposition \ref{prop:kernels-of-powers-increase}(ii) and Corollary \ref{cor:matrix-rank-nullity}(i).\\
The inequality $w_1(\lambda,A)>0$ holds by $\lambda\in\operatorname{spec}(A)$ and Corollary \ref{cor:matrix-rank-nullity}(i).\\
The equality $(\forall p\in\mathbb{Z}_{\geq n+1})\,w_p(\lambda,A)=0$ holds by Theorem \ref{thm:kernel-powers-stabilize-by-dimension} and Corollary \ref{cor:matrix-rank-nullity}(i).\\
Thus $s_{\lambda}(A)\defeq\max\{p\in\{1,\ldots,n\}:w_p(\lambda,A)>0\}$ is well-defined.\\
The \hlt{Weyr characteristic} of $A$ at $\lambda$ is $\bigl(w_p(\lambda,A)\bigr)_{p=1}^{s_{\lambda}(A)}$.
\end{definition}

\begin{theorem} \label{thm:jordan-block-sizes-from-rank-drops} \hlt{Jordan Block Sizes from Rank Drops}\\
Let $n\in\mathbb{Z}_{>0}$, let $A\in\operatorname{Mat}(n,n,\mathbb{C})$, let $\lambda\in\operatorname{spec}(A)$, and let $J=\bigoplus\limits_{\alpha=1}^sJ_{k_{\alpha}}(\mu_{\alpha})$ be a Jordan form of $A$.\\
For $q\in\mathbb{Z}_{\geq0}$, set $r_q\defeq\operatorname{rank}\bigl((A-\lambda I_n)^q\bigr)$; for $q\in\mathbb{Z}_{>0}$, set $w_q\defeq r_{q-1}-r_q$.\\
For $p\in\mathbb{Z}_{>0}$, set $I_p\defeq\{\alpha\in\{1,\ldots,s\}:\mu_{\alpha}=\lambda,\ k_{\alpha}\geq p\}$. Then:
\begin{enumerate}[label=(\roman*)]
\setlength{\itemsep}{0pt}
\item $w_p=\# I_p$;
\item $w_p-w_{p+1}=\#(I_p\setminus I_{p+1})=\#\{\alpha:\mu_{\alpha}=\lambda,\ k_{\alpha}=p\}$.
\end{enumerate}
\begin{proof}
Choose $S\in\operatorname{GL}_n(\mathbb{C})$ such that $A=SJS^{-1}$. The identity $(\forall q\in\mathbb{Z}_{\geq0})\,(A-\lambda I_n)^q=S(J-\lambda I_n)^qS^{-1}$ holds by induction on $q$.\\
Hence $(\forall q\in\mathbb{Z}_{\geq0})\,r_q=\operatorname{rank}((J-\lambda I_n)^q)$ by Theorem \ref{thm:full-rank-products-preserve-rank}(ii). Set $B_{\alpha}\defeq J_{k_{\alpha}}(\mu_{\alpha})$ and $\rho_q(B)\defeq\operatorname{rank}((B-\lambda I)^q)$.\\
Image of $(J-\lambda I_n)^q$ is direct sum of images of $(B_{\alpha}-\lambda I_{k_{\alpha}})^q$ by block-diagonal form. Hence $r_q=\sum\limits_{\alpha=1}^{s}\rho_q(B_{\alpha})$ by Theorem \ref{thm:direct-sum-iff-dimensions-add}(i).\\
Set $\Delta_{p,\alpha}\defeq\rho_{p-1}(B_{\alpha})-\rho_p(B_{\alpha})$. Then $w_p=\sum\limits_{\alpha=1}^{s}\Delta_{p,\alpha}$ by subtracting the two block-rank sums.
\begin{enumerate}[label=(\roman*)]
\setlength{\itemsep}{0pt}
\item Fix $\alpha\in\{1,\ldots,s\}$. If $\mu_{\alpha}\neq\lambda$, then $B_{\alpha}-\lambda I_{k_{\alpha}}=(\mu_{\alpha}-\lambda)I_{k_{\alpha}}+J_{k_{\alpha}}(0)$ is invertible by Lemma \ref{lem:scalar-plus-nilpotent-invertible}.\\
Thus $\Delta_{p,\alpha}=0$ by full rank of every power. If $\mu_{\alpha}=\lambda$, then $\Delta_{p,\alpha}=1$ exactly when $k_{\alpha}\geq p$ by Lemma \ref{lem:nilpotent-jordan-block-power-rank}(i).\\
Thus $\Delta_{p,\alpha}=\mathbf{1}_{I_p}(\alpha)$, where $\mathbf{1}_{I_p}$ is the indicator of $I_p$, and $w_p=\sum\limits_{\alpha=1}^{s}\mathbf{1}_{I_p}(\alpha)=\#I_p$.
\item The equality $w_p-w_{p+1}=\#I_p-\#I_{p+1}=\#(I_p\setminus I_{p+1})$ holds by $I_{p+1}\subseteq I_p$ and part (i).\\
By the definitions of $I_p$ and $I_{p+1}$, $I_p\setminus I_{p+1}=\{\alpha:\mu_{\alpha}=\lambda,\ k_{\alpha}=p\}$.
\end{enumerate}
\hlt{Structural Commentary}: Similarity reduces shifted ranks to block sums; each block contributes a surviving-level indicator.\\
\hlt{Key Insight}: A $\lambda$-block contributes one rank drop through its final level.
\end{proof}
\end{theorem}

Figure \ref{fig:weyr-rank-drops-jordan-blocks} tracks shifted-rank drops to Jordan-block counts.

\begin{illustrationfigure}[H]
  \centering
  \begin{tikzpicture}[
  font=\small,
  rank/.style={draw=blue!60!black, fill=blue!6, rounded corners=2pt, minimum width=0.96cm, minimum height=0.50cm, align=center},
  count/.style={draw=red!60!black, fill=red!6, rounded corners=2pt, minimum width=2.90cm, minimum height=0.50cm, align=center},
  arrow/.style={-{Stealth[length=2mm]}, thick, gray!75}
]
  \node[rank] (r0) at (0,0) {$r_0$};
  \node[rank] (r1) at (1.36,0) {$r_1$};
  \node[rank] (r2) at (2.72,0) {$r_2$};
  \node[rank] (dots) at (4.08,0) {$\cdots$};
  \node[rank] (rq) at (5.44,0) {$r_q$};

  \node[count] (drop) at (2.72,-0.84) {$w_p=r_{p-1}-r_p$};
  \node[count, minimum width=3.28cm] (exact) at (2.72,-1.55) {size-$p$ blocks $=w_p-w_{p+1}$};

  \draw[arrow] (r0) -- (r1);
  \draw[arrow] (r1) -- (r2);
  \draw[arrow] (r2) -- (dots);
  \draw[arrow] (dots) -- (rq);
  \draw[arrow] (drop) -- (exact);
\end{tikzpicture}

  \caption{The $p$th shifted-rank drop counts chains of length at least $p$; consecutive differences count size-$p$ blocks.}
  \label{fig:weyr-rank-drops-jordan-blocks}
\end{illustrationfigure}

\begin{corollary} \label{cor:jordan-form-unique-up-to-block-permutation} \hlt{Uniqueness of Jordan Form}\\
Let $n\in\mathbb{Z}_{>0}$, let $A\in\operatorname{Mat}(n,n,\mathbb{C})$, and let $J,J'$ be Jordan matrices similar to $A$.\\
For $M\in\{J,J'\}$, let $\mathfrak{B}(M)$ be its multiset of Jordan blocks and set $m_M(\lambda,k)\defeq\#\{B\in\mathfrak{B}(M):B=J_k(\lambda)\}$. Then:
\begin{enumerate}[label=(\roman*)]
\setlength{\itemsep}{0pt}
\item $(\forall\,(\lambda,k)\in\operatorname{spec}(A)\times\mathbb{Z}_{>0})\,m_J(\lambda,k)=m_{J'}(\lambda,k)$;
\item $\exists$ a permutation matrix $P$ that permutes whole Jordan-block coordinate groups and satisfies $J'=P^{-1}JP$.
\end{enumerate}
\begin{proof}
\begin{enumerate}[label=(\roman*)]
\setlength{\itemsep}{0pt}
\item Fix $(\lambda,k)\in\operatorname{spec}(A)\times\mathbb{Z}_{>0}$.\\
The equality $m_J(\lambda,k)=w_k(\lambda,A)-w_{k+1}(\lambda,A)=m_{J'}(\lambda,k)$ holds by Theorem \ref{thm:jordan-block-sizes-from-rank-drops}(ii) with $J$ and $J'$ in turn.
\item The ordered block lists of $J$ and $J'$ admit a bijection matching equal blocks by part (i).\\
Let $P$ be the coordinate permutation matrix that carries each block group of $J$ to its matched group in $J'$.\\
Then $P^{-1}JP=J'$ by equality of the matched blocks.
\end{enumerate}
\hlt{Structural Commentary}: Rank-drop multiplicities match equal blocks; a coordinate permutation aligns their ordered groups.\\
\hlt{Key Insight}: Every block multiplicity is intrinsic to the shifted-rank sequence.
\end{proof}
\end{corollary}

\begin{computationmethod} \label{meth:jordan-form-from-rank-drops} \hlt{Jordan Form from Shifted-Rank Drops}\\
Given $n\in\mathbb{Z}_{>0}$ and $A\in\operatorname{Mat}(n,n,\mathbb{C})$, the output is a Jordan matrix $J$ similar to $A$. For $\lambda\in\operatorname{spec}(A)$, perform (i)--(iii).
\begin{enumerate}[label=(\roman*)]
\setlength{\itemsep}{0pt}
\item Compute $r_p(\lambda,A)\defeq\operatorname{rank}((A-\lambda I_n)^p)$ through the first index $t_{\lambda}\in\mathbb{Z}_{>0}$ satisfying $r_{t_{\lambda}-1}(\lambda,A)=r_{t_{\lambda}}(\lambda,A)$.
\item Set $w_p(\lambda,A)\defeq r_{p-1}(\lambda,A)-r_p(\lambda,A)$ for $p\in\{1,\ldots,t_{\lambda}\}$.
\item For $p\in\{1,\ldots,t_{\lambda}-1\}$, include $w_p(\lambda,A)-w_{p+1}(\lambda,A)$ copies of $J_p(\lambda)$.
\item Let $J$ be the direct sum of the blocks obtained over $\lambda\in\operatorname{spec}(A)$.
\end{enumerate}
The stopping index satisfies $t_{\lambda}\leq n+1$ by Theorem \ref{thm:kernel-powers-stabilize-by-dimension} and rank--nullity.\\
The stopping equality implies equality of adjacent shifted kernels by kernel inclusion and rank--nullity.\\
All later shifted-power ranks equal $r_{t_{\lambda}}(\lambda,A)$ by Proposition \ref{prop:kernel-power-stabilization-after-equality} and rank--nullity.\\
The matrix in (iv) is a Jordan form of $A$ by the exact block counts in (iii) and Theorem \ref{thm:jordan-block-sizes-from-rank-drops}(ii).
\end{computationmethod}
